\begin{abstract}
\noindent At the end of a course project on reinforcement learning (a Poisson bandit and two small MDPs) we had three things we could not really explain: how close one can get to the Lai--Robbins bound on a tiny bandit, why TD(0) with step $1/n$ was so slow, and whether results obtained with a single random seed mean anything. This note goes back to each of them on problems where the exact answer is known. On the bandit, KL-UCB and Thompson sampling approach the Lai--Robbins constant only very slowly and from opposite sides, even after $10^6$ pulls, while a UCB with a wrong variance that beats the bound on one instance gets stuck in up to $38\%$ of runs on others. For TD(0), the slowness of $1/n$ is a pure bias effect: with steps $c/n$ the error decays like $n^{-\min(c(1-\gamma),\,1/2)}$, which we see to two decimals, so $1/n$ forgets its initialisation like $n^{-0.2}$ when $\gamma = 0.8$. Finally, rerunning the report on 100 seeds, 7 of the 14 checkable sentences we wrote hold on (almost) every seed, 5 on 80 to 89\% of them, and 2 were mostly an artefact of the seed we happened to use.
\end{abstract}
